\documentclass[12pt]{article}
\usepackage[utf8]{inputenc}
\usepackage{amsmath}
\usepackage{multicol}
\usepackage{geometry}
\usepackage{enumitem}
\usepackage{xcolor}
\usepackage{titlesec}

% Configurações de layout
\geometry{a4paper, left=1cm, right=1cm, top=0.5cm, bottom=1.2cm}
\setlength{\columnseprule}{0.4pt}
\setlength{\baselineskip}{1.0\baselineskip}

% Estilo das seções
\titleformat{\section}
  {\normalfont\Large\bfseries\color{blue}}
  {\thesection}
  {1em}
  {}

\renewcommand{\thesubsection}{\textcolor{red}{\arabic{section}.\arabic{subsection}}}
\titleformat{\subsection}{\color{red}\normalfont\bfseries}{\thesubsection}{1em}{}

\title{\textcolor{blue}{Cultura Digital e Seus Impactos na Sociedade }}
\author{Professor: Jefferson}
\date{}

\begin{document}

\maketitle
 \begin{center}
     \large { \textbf{Leia e responda no caderno}}
    \end{center}
\begin{multicols}{2}
       \section*{Contextualização}
A cultura digital transformou profundamente a maneira como vivemos, nos comunicamos, aprendemos e nos expressamos. Nesta aula, vamos explorar seus conceitos fundamentais, seus impactos na sociedade e os desafios que ela apresenta, com base no artigo da EBDI Corp.

\section*{1. O que é Cultura Digital?}
\subsection*{1.1 Conceito e Características}
Segundo o artigo da EBDI, a \textbf{cultura digital} é o conjunto de práticas, comportamentos, valores, crenças e expressões que surgem da interação das pessoas com as tecnologias digitais e a Internet.

\textbf{Principais características:}
\begin{itemize}
    \item \textbf{Participação ativa}: As pessoas não apenas consomem conteúdo digital, mas também o criam, compartilham e modificam.
    \item \textbf{Ambiente dinâmico}: Plataformas como mídias sociais, blogs e fóruns estão em constante evolução.
    \item \textbf{Democratização do acesso}: Informação disponível a qualquer momento e lugar, impactando educação, pesquisa e jornalismo.
    \item \textbf{Expressão criativa}: Arte, música, vídeos e literatura alcançam públicos globais de forma colaborativa.
\end{itemize}

\subsection*{1.2 Exemplos Práticos}
\begin{itemize}
    \item \textbf{Criação de conteúdo}: YouTubers, podcasters, criadores de memes e artistas independentes.
    \item \textbf{Plataformas colaborativas}: Wikipedia, GitHub, fóruns especializados.
    \item \textbf{Educação online}: Cursos gratuitos, tutoriais, aplicativos de idiomas.
\end{itemize}

\section*{2. Impactos da Cultura Digital na Sociedade}
O artigo destaca vários papéis fundamentais da cultura digital. Vamos analisar cada um:

\subsection*{2.1 Acesso à Informação}
A Internet ampliou enormemente o acesso à informação. Hoje, é possível obter conhecimento em praticamente qualquer campo, a qualquer momento. Isso transformou a educação, a pesquisa e o jornalismo, permitindo o aprendizado contínuo e a troca de ideias em escala global.

\subsection*{2.2 Transformação da Comunicação}
Plataformas digitais proporcionam meios instantâneos e globais de interação. Redes sociais, aplicativos de mensagens e videoconferências conectam pessoas em tempo real, independentemente da distância geográfica.

\subsection*{2.3 Criação e Expressão Cultural}
A cultura digital permite que pessoas criem, compartilhem e consumam diversas formas de arte. Artistas independentes alcançam públicos globais sem a necessidade de grandes orçamentos ou intermediários tradicionais.

\subsection*{2.4 Ativismo e Mobilização Social}
As mídias sociais desempenham papel crucial na organização de movimentos sociais, protestos e campanhas de conscientização. Elas fornecem espaços para expressão política, mobilização em torno de causas e denúncia de injustiças.

\subsection*{2.5 Economia Digital}
A cultura digital impulsiona novos negócios: startups de tecnologia, comércio eletrônico, freelances digitais. Setores tradicionais como entretenimento e publicidade foram transformados pela digitalização.

\subsection*{2.6 Educação}
A tecnologia digital está mudando fundamentalmente a maneira como aprendemos e ensinamos. Plataformas de aprendizado online, tutoriais em vídeo, aplicativos educacionais e jogos educacionais são ferramentas que remodelam a educação.

\section*{3. Desafios da Cultura Digital}
Apesar dos benefícios, o artigo destaca preocupações importantes:

\subsection*{3.1 Privacidade e Segurança Online}
A coleta e uso de dados pessoais por plataformas digitais levanta questões sobre privacidade e vigilância.

\subsection*{3.2 Desigualdade de Acesso}
O acesso à tecnologia não é universal. Diferenças socioeconômicas criam exclusão digital, limitando oportunidades.

\subsection*{3.3 Disseminação de Desinformação}
A facilidade de compartilhar informações também permite a propagação de notícias falsas e teorias conspiratórias.

\subsection*{3.4 Necessidade de Abordagem Ética}
Todas essas questões exigem uma reflexão responsável e ética sobre o uso das tecnologias digitais.

\section*{4. Atividade}
\subsection*{Questão 1 – Questões Conceituais}
Responda em seu caderno:
\begin{enumerate}[label=\alph*)]
    \item Defina cultura digital com suas próprias palavras, utilizando pelo menos três características mencionadas no texto.
    \item Explique a diferença entre consumo passivo e participação ativa no ambiente digital. Dê um exemplo de cada.
    \item Por que o artigo afirma que a cultura digital é um “fenômeno complexo e multifacetado”?
\end{enumerate}

\subsection*{Questão 2 – Análise de Impactos}
\begin{enumerate}[label=\alph*)]
    \item Escolha duas áreas impactadas pela cultura digital (comunicação, educação, economia, política) e descreva uma mudança significativa em cada uma.
    \item Como a democratização do acesso à informação transformou a educação e o jornalismo? Cite exemplos concretos.
    \item De que forma a cultura digital permite que artistas independentes alcancem públicos globais sem intermediários tradicionais?
\end{enumerate}

\subsection*{Questão 3 – Desafios Éticos e Sociais}
\begin{enumerate}[label=\alph*)]
    \item Explique por que a desigualdade de acesso à tecnologia é considerada um problema social. Quais grupos são mais afetados?
    \item Como a disseminação de desinformação pode afetar a democracia e a tomada de decisões?
    \item O que você entende por “abordagem responsável e ética” no uso das tecnologias digitais? Dê um exemplo.
\end{enumerate}

\subsection*{Questão 4 – Pesquisa e Reflexão}
\begin{enumerate}[label=\alph*)]
    \item Descreva um movimento social ou campanha de conscientização que tenha utilizado redes sociais como ferramenta de organização. Descreva como isso aconteceu.
    \item Escolha uma plataforma digital (YouTube, TikTok, Instagram, etc.) e analise: que tipo de participação ativa ela permite? Quais são os benefícios e os riscos associados?
    \item Escreva um pequeno texto (10 a 15 linhas) respondendo: \textit{"Como a cultura digital impacta minha vida pessoal e minha forma de aprender?"}
\end{enumerate}

\end{multicols}

\end{document}
